\section{长上下文下的注意力计算}
\label{sec:nlp-llm-long-context}

第\ref{sec:nlp-llm-computation}节的解码器让每个位置按式\eqref{eq:nlp-llm-causal-attention}读取全部前缀。上下文从几千词元增加到十几万词元时，这一读取方式同时暴露出三个问题：训练时注意力得分的计算随长度平方增长；推断时每生成一个词元都要保存并读取全部历史的键值状态；而模型能否精确回读任意早先位置，又决定了查找、复制一类任务能做到什么程度。注意力、多头、键值缓存、MQA与GQA以及线性注意力的通用定义见第\ref{sec:nlp-llm-computation}节开头所引卷五的Transformer章；本节只写大语言模型中的配置、推导与比较。

按改变的对象，长上下文方法可分为四类。第一类保留读取全部历史，只压缩每个位置保存的键值，包括MQA、GQA与MLA；第二类保留每个位置的完整键值，但每次只读取其中一部分，包括滑动窗口、全局词元与NSA；第三类不再逐项保存历史，而把历史汇聚为固定大小的递归状态，包括线性注意力、RetNet与GLA、delta规则与Gated DeltaNet、Mamba；第四类把递归层与少量全注意力层混合。FlashAttention只改变访存顺序而不改变数学结果和计算量，放在第\ref{chap:nlp-model-inference}章讨论。

\subsection{计算量与状态量的瓶颈}
\label{sec:nlp-llm-context-bottlenecks}

先用第\ref{sec:nlp-llm-computation}节的LLaMA-3-8B配置估算全注意力的代价：层数$L=32$，宽度$d=4096$，查询头数$H=32$，键值头数$H_{\mathrm{kv}}=8$，每头维度$d_h=128$，参数量$N\approx8.0\times10^9$。

训练计算方面，式\eqref{eq:nlp-llm-forward-flops}的第二项$2LTd$就是注意力得分与加权求和这两步不含参数的乘法，已按因果掩码取平均。它对每个词元随$T$线性增长，对整段序列则随$T^2$增长。$T=8192$时这一项约为$2.1\times10^9$，只占参数项$2N_{\mathrm{mat}}\approx1.5\times10^{10}$的14\%；$T=131072$时增长到约$3.4\times10^{10}$，是参数项的2.3倍。长度扩大16倍，整段序列的得分计算扩大约$16^2=256$倍。

推断状态方面，生成时要为每层每个历史词元保存键和值，每个词元的元素数为$2LH_{\mathrm{kv}}d_h=2\times32\times8\times128=65536$。以每个元素2字节的bf16计，每个词元约占128\,KiB，13.1万词元的上下文就需要约16\,GiB缓存，与80亿参数本身的bf16存储量相当。更关键的是读取量：逐步解码时每生成一个词元，都要把整份缓存从显存读一遍，序列越长，每步越受显存带宽限制。

第三个维度是信息访问。全注意力在任意一步都能以相同方式读到任意一个历史位置，因此可以从很长的前缀中精确复制一个编号或一处改动。后文的方法分别在三处取舍：减少每个位置的状态，减少每次读取的位置，或者放弃逐位置保存。训练计算、推断状态与读取量、精确回读能力很难同时取到最优，我们在第\ref{sec:nlp-llm-context-comparison}节用一张表把这三项放在一起比较。

\subsection{键值状态的压缩}
\label{sec:nlp-llm-kv-compression}

第一类方法不改变“每个查询读取全部历史”的规则，只减少每个历史位置需要保存的键值量。MQA与GQA减少键值头的个数；MLA保留多头的读取方式，但让所有头的键值都从一个低维潜在向量生成，并借助矩阵吸收避免在推断时把它展开。两者都缩小推断缓存，训练时得分矩阵的平方计算不变。

\subsubsection{多查询注意力与分组查询注意力}
\label{sec:nlp-llm-mqa-gqa}

多头注意力中每个查询头有自己的键头和值头，缓存量与$H$成正比。\term{多查询注意力}（multi-query attention, MQA）让所有查询头共享一组键和值，每层每个词元只存$2d_h$个元素\scite{shazeer2019mqa}。第\ref{sec:nlp-llm-causal-mha}小节的GQA是两者之间的折中：$H$个查询头分成$H_{\mathrm{kv}}$组，组内共享一组键值，缓存量为$2H_{\mathrm{kv}}d_h$，是多头注意力的$H_{\mathrm{kv}}/H$倍\scite{ainslie2023gqa}。$H_{\mathrm{kv}}=H$时它退化为多头注意力，$H_{\mathrm{kv}}=1$时就是MQA。LLaMA-3-8B取$H=32$、$H_{\mathrm{kv}}=8$，缓存缩为四分之一。

GQA的计算与多头注意力只差一处索引，即第\ref{sec:nlp-llm-causal-mha}小节的组映射$g(a)$。实现时可以把键值张量从形状$(T,H_{\mathrm{kv}},d_h)$沿头维复制成$(T,H,d_h)$再照常计算，或直接在核函数中按组索引读取。已有的多头检查点也可以转成GQA：把每组内各头的键投影矩阵和值投影矩阵分别取平均，作为该组共享的投影，再用原预训练计算量的一小部分继续训练，让其余参数适应新的共享方式；Ainslie等报告约5\%的继续训练即可接近原模型质量\scite{ainslie2023gqa}。共享的代价是同组查询头只能在同一组键值上做匹配，键值表示的自由度随$H_{\mathrm{kv}}$减小而下降；MQA在大模型上常见的质量损失正来自这里。

\subsubsection{多头潜在注意力}
\label{sec:nlp-llm-mla}

GQA通过减少头数节省缓存。\term{多头潜在注意力}（multi-head latent attention, MLA）换一种约束：每个头仍有自己的键和值，但它们全部由同一个低维潜在向量线性生成，缓存只保存这个潜在向量\scite{nlp-llm-deepseekv2-2024}。下面先省略位置编码，写出内容分支：
\begin{equation}
 \bm c_j=\bm W_{D}\bm h_j,\qquad
 \bm k_{j,i}^{C}=\bm W_{UK,i}\bm c_j,\qquad
 \bm v_{j,i}^{C}=\bm W_{UV,i}\bm c_j,
 \label{eq:nlp-llm-mla}
\end{equation}
其中$\bm h_j\in\mathbb R^d$是位置$j$的层输入，$\bm c_j\in\mathbb R^{d_c}$是共享的潜在向量，$d_c$远小于展开后的$2Hd_h$。下投影$\bm W_D\in\mathbb R^{d_c\times d}$为所有头共用；$\bm W_{UK,i}\in\mathbb R^{d_h\times d_c}$与$\bm W_{UV,i}\in\mathbb R^{d_h\times d_c}$是第$i$个头的上投影，把它们按头纵向堆叠就得到形状为$Hd_h\times d_c$的完整上投影矩阵。

如果推断时把$\bm c_j$展开成各头键值再计算，缓存虽然变小，计算却要为每个历史位置重做上投影。MLA利用线性性避开这一步。设第$i$个头的查询为$\bm q_{t,i}=\bm W_{Q,i}\bm h_t$，则
\[
 \bm q_{t,i}^{\top}\bm k_{j,i}^{C}
 =\bm q_{t,i}^{\top}\bm W_{UK,i}\bm c_j
 =\bigl(\bm W_{UK,i}^{\top}\bm W_{Q,i}\bm h_t\bigr)^{\top}\bm c_j .
\]
括号内只依赖当前位置，矩阵$\bm W_{UK,i}^{\top}\bm W_{Q,i}\in\mathbb R^{d_c\times d}$可以预先乘好，于是每个头直接在$d_c$维潜在空间中与历史的$\bm c_j$做内积。值的聚合同样可以交换次序：
\begin{equation}
 \bm o_{t,i}=\sum_{j=1}^{t}a_{t,i,j}\bm W_{UV,i}\bm c_j
 =\bm W_{UV,i}\sum_{j=1}^{t}a_{t,i,j}\bm c_j,
 \label{eq:nlp-llm-mla-value}
\end{equation}
先在潜在向量上加权求和，再做一次上投影；$\bm W_{UV,i}$还可以并入输出投影$\bm W_O$中对应第$i$头的列块。这一做法称为投影吸收：上投影矩阵被吸收进查询侧与输出侧，历史缓存始终只有$\bm c_j$。代价是每个得分的内积维度从$d_h$变为$d_c$，单步计算量上升；解码阶段主要受显存读取限制，用计算换读取通常是划算的。

位置编码打破了上述吸收。若对查询和键施加RoPE，得分变为
\[
 (\bm R_t\bm q_{t,i})^{\top}\bm R_j\bm W_{UK,i}\bm c_j=\bm q_{t,i}^{\top}\bm R_{j-t}\bm W_{UK,i}\bm c_j ,
\]
旋转矩阵$\bm R_{j-t}$夹在查询与$\bm W_{UK,i}$之间且随历史位置$j$变化，不能再预先把$\bm W_{UK,i}$并入查询。DeepSeek-V2因此把位置信息放在一条单独的解耦分支上：每个头另有承载RoPE的位置查询$\bm q_{t,i}^{R}=\bm R_t\bm W_{QR,i}\bm h_t\in\mathbb R^{d_R}$，所有头共享一个位置键$\bm k_j^{R}=\bm R_j\bm W_{KR}\bm h_j\in\mathbb R^{d_R}$，得分写为
\[
 s_{t,i,j}=\frac{(\bm q_{t,i}^{C})^{\top}\bm k_{j,i}^{C}+(\bm q_{t,i}^{R})^{\top}\bm k_j^{R}}{\sqrt{d_h+d_R}} .
\]
内容部分照常吸收，位置部分维度很小且各头共享，历史缓存为$\bm c_j$与$\bm k_j^{R}$两项。DeepSeek-V2还对查询做类似的低秩压缩，那只减少训练时的激活存储，与缓存无关。

用一个教学配置比较每层每个历史词元的缓存元素数：取$H=32$、$d_h=128$。多头注意力保存$2Hd_h=8192$个元素；$H_{\mathrm{kv}}=8$的GQA保存$2H_{\mathrm{kv}}d_h=2048$个；取$d_c=4d_h=512$、$d_R=d_h/2=64$的MLA保存$d_c+d_R=576$个，相当于约2.25组键值头的GQA。DeepSeek-V2正是取$d_c=4d_h$、$d_R=d_h/2$。这一比较只统计状态元素，整体运行时间还取决于吸收后的计算量与具体实现，不会按同一比例变化。GQA与MLA的区别可以概括为：前者限制不同头使用的键值必须相同，后者限制所有头的键值落在同一个$d_c$维子空间中，后一种约束允许各头仍有不同的键值。

\subsection{历史位置的稀疏读取}
\label{sec:nlp-llm-sparse-reading}

键值压缩没有减少查询与历史位置的配对数，训练时得分计算仍是$O(T^2)$。第二类方法保留每个位置的完整键值，但规定每个查询只读取其中一部分，这就是\term{稀疏注意力}（sparse attention）。读取范围可以按固定模式规定，也可以按内容选择。固定模式实现简单、计算可预知，远距离信息只能经多层间接传递；按内容选择能直接访问远处的相关块，但选择本身需要计算，也可能选漏。

\subsubsection{滑动窗口与全局词元}
\label{sec:nlp-llm-sliding-window}

\term{滑动窗口注意力}（sliding window attention）让位置$t$只读取最近的$w$个位置$\{t-w+1,\ldots,t\}$，每层的得分计算从$O(T^2)$降为$O(Tw)$。一层只能看到窗口内，但第$\ell$层位置$t$的表示已经汇总了第$\ell-1$层$[t-w+1,t]$中各位置的信息，而这些位置又各自汇总了更早的$w$个位置，因此$L$层叠加后的感受野约为$Lw$。以$w=4096$、$L=32$计，理论上能触及约13万词元之前的位置。远处信息要逐层接力传递，每传一层都与其他内容一起被压缩进有限维表示，所以可达不等于可精确读取。推断时只需保留每层最近$w$个位置的键值，Mistral 7B就用这种滚动缓存把缓存上限固定在窗口长度\scite{nlp-llm-mistral2023}。

纯窗口缺少汇总全局信息的位置。Longformer在窗口之外为少数指定位置设置\term{全局词元}（global token）：它们读取所有位置，也被所有位置读取\scite{nlp-llm-longformer2020}。在问答中可以把问题词元设为全局，让每个正文位置都能直接对照问题，问题位置也能汇总全文。全局词元数$g$固定时，每层代价为$O(T(w+g))$。在因果语言模型中，“被所有位置读取”只能对排在前面的全局词元成立，常见做法是把开头若干词元当作始终可见的位置。大语言模型的注意力权重往往大量落在序列最前面的几个词元上，这一注意力汇点现象及据此设计的StreamingLLM缓存策略放在第\ref{chap:nlp-model-inference}章讨论。

窗口和全局词元都按位置而不是按内容决定读取范围。若“读书会改到周五”这条通知在窗口之外，当前位置只能指望它已被逐层接力的表示带过来，或恰好写进了全局位置。

\subsubsection{原生稀疏注意力}
\label{sec:nlp-llm-nsa}

\term{原生稀疏注意力}（native sparse attention, NSA）让当前查询按内容决定读取远处的哪些块，并从预训练开始就使用这一稀疏结构，而不是在训练好的稠密模型上事后剪枝\scite{nlp-llm-nsa2025}。它设三个并行分支，各自对当前查询$\bm q_t$计算一次注意力，再用门控融合：
\begin{equation}
 \bm o_t=g_t^{c}\bm o_t^{c}+g_t^{s}\bm o_t^{s}+g_t^{w}\bm o_t^{w},
 \label{eq:nlp-llm-nsa}
\end{equation}
上标$c,s,w$分别表示压缩、选择与滑动窗口分支，门值$g_t^{b}\in[0,1]$由当前表示经小型网络与sigmoid得到。三个分支分别归一化，不是把三类键值拼接后共用一次softmax。

压缩分支提供粗粒度的全局视图。它把历史键按长度$l$、步长$\delta$切成连续块，每块经一个带块内位置编码的小型网络$\varphi$压成一个块级键，值也同样处理。位置$t$可读取的压缩键约有$t/\delta$个，查询对它们做普通注意力，得到$\bm o_t^{c}$和一组块级权重$\bm p_t^{c}$。

选择分支用压缩分支的结果决定读取哪些原始块。历史被另切成长度为$l'$的选择块；一个选择块的重要性分数取与它重叠的各压缩块权重之和，因此选块几乎不增加额外的打分计算。取分数最高的$n$个选择块，查询读取这些块中全部原始词元的键值，得到$\bm o_t^{s}$。为了配合GQA，同一组内各查询头的块分数先相加再统一选块，使一组头读取同一批连续的键值，这样能按块连续地从显存加载数据。滑动窗口分支只读取最近$w$个位置，承担局部依赖，避免局部模式在训练中“抢走”另两个分支的学习信号；原文还为三个分支分别配置键值投影，进一步减少分支之间的职责干扰。

每步读取的位置数约为$t/\delta+nl'+w$。按原文配置$l=32$、$\delta=16$、$l'=64$、$n=16$、$w=512$，在$t=65536$时约为$4096+1024+512=5632$，约为全注意力的十二分之一。选块是离散的top-$n$操作，梯度只经过被选中块的注意力和门控传回，压缩分支的参数通过其自身输出$\bm o_t^{c}$得到训练。

选漏的后果需要单独说明。假设很长的活动记录中，地点曾在早先某一块里被改为“二楼会议室”，当前问题询问最终地点。若这一块的压缩表示与当前查询匹配不高，没有进入前$n$块，选择分支就读不到原文；压缩分支只能提供被压缩过的粗略信息，局部窗口也够不着，后续计算再准确也无法恢复这条证据。NSA的收益在于读取量随长度近似线性增长且保留按内容访问远处原文的能力；代价是选择质量依赖压缩表示，所有键值仍需保存，缓存量并不减少。

\subsection{历史信息的递归汇聚}
\label{sec:nlp-llm-recurrent-state}

前两类方法都保留了每个历史位置的键值。第三类方法把历史写入一个大小固定的状态矩阵，每步推断只需常数大小的状态和常数计算，训练时又可以分块并行。代价是有限状态必然有损：它无法像全注意力那样精确回读任意早先位置。各方法的区别集中在状态怎样写入与遗忘：线性注意力只累加；RetNet与GLA加入固定或数据相关的衰减；delta规则按当前键修正旧关联；Mamba以输入相关的离散化参数控制状态演化。图~\ref{fig:nlp-llm-context-state-update}对照了全注意力的键值缓存与递归状态的更新方式。

\begin{figure*}[!t]
  \centering
  % 第九卷第3章：全注意力的键值缓存与递归汇聚的状态更新对照（跨栏，约150 mm宽）。
\begin{tikzpicture}[
  x=1mm,y=1mm,
  figure picture,
  cell/.style={rectangle,draw=FigureInput,fill=FigureInputFill,line width=.6pt,
    minimum width=10mm,minimum height=6mm,inner sep=0pt,font=\figurefont\footnotesize},
  vcell/.style={cell,draw=FigureModel,fill=FigureModelFill},
  qnode/.style={figure node,draw=FigureOutput,fill=FigureOutputFill,minimum height=7mm,
    minimum width=10mm},
  state/.style={figure node,draw=FigureModel,fill=FigureModelFill,minimum width=16mm,
    minimum height=14mm},
  op/.style={figure node,minimum height=9mm,minimum width=26mm},
  dop/.style={op,dash pattern=on 2pt off 1.5pt},
  plus/.style={circle,draw=FigureStroke,line width=.75pt,minimum size=5mm,inner sep=0pt,
    font=\figurefont\footnotesize},
  lab/.style={figure note,anchor=west}
]
  % ---------- 上方：全注意力的键值缓存 ----------
  \node[figure heading,anchor=west] at (0,30) {全注意力：逐项追加键值缓存};
  \foreach \i/\name in {0/1,1/2,2/3} {
    \node[cell] (k\name) at ({8+\i*10},20) {$\bm k_{\name}$};
    \node[vcell] (v\name) at ({8+\i*10},13) {$\bm v_{\name}$};
  }
  \node[font=\figurefont\footnotesize] at (38,16.5) {$\cdots$};
  \node[cell] (kt1) at (48,20) {$\bm k_{t-1}$};
  \node[vcell] (vt1) at (48,13) {$\bm v_{t-1}$};
  \node[cell,dash pattern=on 2pt off 1.5pt] (kt) at (58,20) {$\bm k_{t}$};
  \node[vcell,dash pattern=on 2pt off 1.5pt] (vt) at (58,13) {$\bm v_{t}$};
  \draw[figure link] (3,8) -- (63,8);
  \draw[figure link] (3,7) -- (3,9);
  \draw[figure link] (63,7) -- (63,9);
  \node[figure note] at (33,4.5) {缓存长度$t$，随生成线性增长};
  \node[qnode] (q) at (82,16.5) {$\bm q_t$};
  \foreach \n in {k1,k2,k3,kt1,kt} {
    \draw[figure link,densely dotted] (\n.north) -- (\n.north |- 0,26);
  }
  \draw[figure link,densely dotted] (k1.north |- 0,26) -- (70,26) |- (q.west);
  \node[lab] at (90,19) {与全部$t$个键配对打分，};
  \node[lab] at (90,14) {再对$t$个值加权求和};

  % 分隔线
  \draw[draw=FigureGrid,line width=.5pt] (0,0) -- (150,0);

  % ---------- 下方：递归汇聚的固定大小状态 ----------
  \node[figure heading,anchor=west] at (0,-6) {递归汇聚：更新固定大小的状态矩阵};
  \node[state] (sold) at (10,-24) {$\bm S_{t-1}$\\[-1pt]{\scriptsize $d_v\times d_k$}};
  \node[dop] (decay) at (45,-24) {衰减或擦除\\{\scriptsize 修改旧状态}};
  \node[plus] (add) at (72,-24) {$+$};
  \node[op,draw=FigureInput,fill=FigureInputFill] (write) at (72,-42) {写入外积\\$\beta_t\bm v_t\bm k_t^{\top}$};
  \node[state] (snew) at (98,-24) {$\bm S_{t}$\\[-1pt]{\scriptsize $d_v\times d_k$}};
  \node[qnode] (q2) at (98,-45) {$\bm q_t$};
  \node[figure node,draw=FigureOutput,fill=FigureOutputFill,minimum height=9mm,minimum width=22mm] (out) at (132,-24) {读出\\$\bm o_t=\bm S_t\bm q_t$};

  \draw[figure flow] (sold.east) -- (decay.west);
  \draw[figure flow] (decay.east) -- (add.west);
  \draw[figure flow] (write.north) -- (add.south);
  \draw[figure flow] (add.east) -- (snew.west);
  \draw[figure flow] (snew.east) -- (out.west);
  \draw[figure flow] (q2.east) -| (out.south);
  \draw[figure flow,draw=FigureMuted,dash pattern=on 2pt off 1.5pt]
    (snew.north) -- ++(0,6) -| (sold.north);
  \node[figure note] at (54,-13.5) {下一步作为旧状态};
\end{tikzpicture}

  \caption{全注意力与递归汇聚保存历史的方式。上方：全注意力把每个历史位置的键$\bm k_j$和值$\bm v_j$逐项追加到缓存，当前查询$\bm q_t$与全部$t$个键配对，缓存随$t$线性增长。下方：递归汇聚只保留一个$d_v\times d_k$的状态矩阵$\bm S$，每步先按衰减或擦除修改旧状态，再写入外积$\bm v_t\bm k_t^{\top}$，最后以$\bm o_t=\bm S_t\bm q_t$读出；状态大小与$t$无关。虚线框中的“衰减或擦除”在线性注意力中为恒等，在RetNet中为$\gamma\bm S$，在GLA中为$\bm S\operatorname{Diag}(\bm\alpha_t)$，在Gated DeltaNet中为$\alpha_t\bm S(\bm I-\beta_t\bm k_t\bm k_t^{\top})$，写入项相应乘以$\beta_t$。}
  \label{fig:nlp-llm-context-state-update}
\end{figure*}

\subsubsection{线性注意力及其递推形式}
\label{sec:nlp-llm-linear-attention}

softmax注意力的难点在于$\exp(\bm q_t^{\top}\bm k_j)$不能拆成只依赖$\bm q_t$与只依赖$\bm k_j$的两部分之积，所以必须对每对位置单独计算。\term{线性注意力}（linear attention）把这个相似度换成特征映射的内积$\phi(\bm q_t)^{\top}\phi(\bm k_j)$，其中$\phi:\mathbb R^{d_h}\to\mathbb R^{r}$输出非负向量，例如Katharopoulos等取$\phi(\bm x)=\operatorname{elu}(\bm x)+1$\scite{katharopoulos2020linear}。代入因果注意力，矩阵乘法的结合律允许先对历史求和：
\[
 \bm o_t=\frac{\sum_{j\le t}\bm v_j\phi(\bm k_j)^{\top}\phi(\bm q_t)}{\sum_{j\le t}\phi(\bm k_j)^{\top}\phi(\bm q_t)}
 =\frac{\bm S_t\phi(\bm q_t)}{\bm z_t^{\top}\phi(\bm q_t)},
\]
其中
\[
 \bm S_t=\sum_{j\le t}\bm v_j\phi(\bm k_j)^{\top},\qquad
 \bm z_t=\sum_{j\le t}\phi(\bm k_j).
\]
$\bm S_t\in\mathbb R^{d_v\times r}$和$\bm z_t\in\mathbb R^{r}$可以逐步累加。后续改进普遍去掉分母，并直接把经过归一化的$\bm k_t,\bm q_t$当作特征（$r=d_k$），得到最简洁的递推形式：
\begin{equation}
 \bm S_t=\bm S_{t-1}+\bm v_t\bm k_t^{\top},\qquad
 \bm o_t=\bm S_t\bm q_t .
 \label{eq:nlp-llm-linear-recurrence}
\end{equation}

$\bm S_t$可以理解为一个\term{关联记忆}（associative memory）：每写入一对$(\bm k_t,\bm v_t)$，就在矩阵中叠加一个外积；用查询读取时，$\bm S_t\bm q_t=\sum_j\bm v_j(\bm k_j^{\top}\bm q_t)$，与查询方向相近的键所对应的值得到较大权重。这也说明它与前馈网络中“权重存储知识”的关系：$\bm S_t$是随输入变化的一组快速权重，Schlag等因此把线性注意力称为快速权重编程器\scite{nlp-llm-fastweight2021}。推断时每步只需更新和读取一个$d_v\times d_k$矩阵，单头计算为$O(d_kd_v)$，与上下文长度无关。

训练时逐步递推无法利用矩阵乘法的并行度，常用分块形式在并行与递推之间折中。把序列切成长为$C$的块，块内用带因果掩码的矩阵乘法并行计算，块间只传递状态矩阵。设第$c$块的查询、键、值按行堆叠为$\bm Q_c,\bm K_c\in\mathbb R^{C\times d_k}$、$\bm V_c\in\mathbb R^{C\times d_v}$，$\bm M$为含对角线的下三角全1矩阵，计算过程见算法~\ref{alg:nlp-llm-chunk-linear}。

\begin{algorithm}[线性注意力的分块并行计算]
  \label{alg:nlp-llm-chunk-linear}
  \begin{algorithmic}[1]
    \Require 按块切分的$\bm Q_c,\bm K_c,\bm V_c$，$c=1,\ldots,T/C$；块内因果掩码$\bm M\in\{0,1\}^{C\times C}$
    \Ensure 各块输出$\bm O_c\in\mathbb R^{C\times d_v}$
    \State $\bm S\gets\bm 0\in\mathbb R^{d_v\times d_k}$
    \For{$c=1,\ldots,T/C$}
      \State $\bm O_c\gets\bm Q_c\bm S^{\top}+\bigl((\bm Q_c\bm K_c^{\top})\odot\bm M\bigr)\bm V_c$ \Comment{块间读状态，块内做因果注意力}
      \State $\bm S\gets\bm S+\bm V_c^{\top}\bm K_c$ \Comment{把本块写入状态}
    \EndFor
  \end{algorithmic}
\end{algorithm}

第一项读取此前各块累积的状态，第二项处理块内配对，两项相加恰好等于式\eqref{eq:nlp-llm-linear-recurrence}逐步展开的结果。整段计算量约为$O(TCd_k+Td_kd_v)$，取$C$为几十到几百即可让矩阵乘法充分利用硬件，而不必形成$T\times T$的得分矩阵。

线性注意力的缺点也来自式\eqref{eq:nlp-llm-linear-recurrence}。状态只加不减，旧关联永远留在矩阵中；$d_k$维空间最多容纳$d_k$个相互正交的键，写入更多键后，读取任意一个键都会混入其他值的成分。没有遗忘机制时，这种干扰随长度累积，早期线性注意力在语言建模上明显落后于softmax注意力，主要原因就在这里。

\subsubsection{衰减与门控：RetNet与GLA}
\label{sec:nlp-llm-decay-gating}

最直接的改进是让旧记忆随时间衰减。\term{保持网络}（retentive network, RetNet）在每个头上使用固定衰减率$\gamma\in(0,1)$\scite{nlp-llm-retnet2023}：
\begin{equation}
 \bm S_t=\gamma\bm S_{t-1}+\bm v_t\bm k_t^{\top},\qquad \bm o_t=\bm S_t\bm q_t .
 \label{eq:nlp-llm-retnet}
\end{equation}
展开得$\bm o_t=\sum_{j\le t}\gamma^{t-j}(\bm k_j^{\top}\bm q_t)\bm v_j$，即距离为$t-j$的历史按$\gamma^{t-j}$打折。这一形式有三种等价算法：并行形式写成$\bm O=\bigl((\bm Q\bm K^{\top})\odot\bm D\bigr)\bm V$，其中$D_{t,j}=\gamma^{t-j}$（$t\ge j$）否则为0，适合训练；递归形式即式\eqref{eq:nlp-llm-retnet}，适合推断；分块形式在块内用并行形式、块间传递状态，衰减因子按块内位置补乘。RetNet给不同头设不同的$\gamma$，使有的头关注近处、有的头保留较远的历史；查询与键上还施加与位置相关的旋转，作用与RoPE类似。

固定衰减对所有内容一视同仁：无论当前词元是否重要，旧状态都按同样比例缩小。\term{门控线性注意力}（gated linear attention, GLA）改用数据相关的逐维衰减\scite{nlp-llm-gla2024}：
\begin{equation}
 \bm S_t=\bm S_{t-1}\operatorname{Diag}(\bm\alpha_t)+\bm v_t\bm k_t^{\top},\qquad
 \bm\alpha_t=\sigma(\bm W_{\alpha}\bm h_t+\bm b_{\alpha})^{1/\tau}\in(0,1)^{d_k},
 \label{eq:nlp-llm-gla}
\end{equation}
右乘对角矩阵按键维度逐列缩放状态，$\tau>1$是使门值接近1的温度，$\bm W_{\alpha}$常取低秩分解以节省参数。当前输入若示意“话题已经切换”，门值可以变小，较快地清除旧内容；若是无关的填充词元，门值接近1，旧内容得以保留。门控随位置变化后，并行形式中的衰减矩阵不再是$\gamma^{t-j}$，而是各步门值的累积乘积；GLA把这些乘积并入块内计算，仍然保持分块并行训练。衰减与门控缓解了干扰的累积，但它们只能整体缩小旧内容，不能有针对性地改写某个键下的旧值。

\subsubsection{delta规则与Gated DeltaNet}
\label{sec:nlp-llm-delta-rule}

考虑“读书会改到周五”这类更正：同一个键（活动时间）先后对应不同的值。线性注意力会把新旧两个值叠加在一起，衰减也只能让二者一起变弱。要只改这一个键下的值，可以把状态的更新看作一个在线回归问题：希望状态把当前键映射到当前值，即最小化$\ell_t(\bm S)=\frac12\lVert\bm S\bm k_t-\bm v_t\rVert_2^2$。从$\bm S_{t-1}$出发、以步长$\beta_t$做一步梯度下降，梯度为$(\bm S_{t-1}\bm k_t-\bm v_t)\bm k_t^{\top}$，于是
\begin{equation}
 \bm S_t=\bm S_{t-1}-\beta_t(\bm S_{t-1}\bm k_t-\bm v_t)\bm k_t^{\top}
 =\bm S_{t-1}(\bm I-\beta_t\bm k_t\bm k_t^{\top})+\beta_t\bm v_t\bm k_t^{\top}.
 \label{eq:nlp-llm-delta}
\end{equation}
这就是\term{delta规则}（delta rule）：$\bm S_{t-1}\bm k_t$是状态对当前键的“旧答案”，括号里的误差决定修正多少，修正只沿$\bm k_t$方向进行\scite{nlp-llm-fastweight2021}。$\beta_t\in(0,1)$由当前输入经sigmoid产生，表示写入强度。第二种写法给出另一种读法：先用$\bm I-\beta_t\bm k_t\bm k_t^{\top}$擦除状态在$\bm k_t$方向上的旧内容，再写入新值。线性注意力式\eqref{eq:nlp-llm-linear-recurrence}对应的在线目标是$-\langle\bm S\bm k_t,\bm v_t\rangle$，梯度与状态无关，所以只能累加。

用一个二维小例子比较两种规则（数值为教学示例）。键取单位向量$\bm k=(1,0)^{\top}$，先写入值$\bm a=(2,0)^{\top}$，得$\bm S=\bm a\bm k^{\top}$；随后同一键对应新值$\bm b=(0,3)^{\top}$。线性注意力得$\bm S'=(\bm a+\bm b)\bm k^{\top}$，用$\bm k$读出$(2,3)^{\top}$，新旧值混在一起。delta规则取$\beta=1$时$\bm S'=\bm S(\bm I-\bm k\bm k^{\top})+\bm b\bm k^{\top}=\bm b\bm k^{\top}$，读出恰为$(0,3)^{\top}$；与$\bm k$正交的另一个键$(0,1)^{\top}$下存的值不受影响。一般地，键经过$L_2$归一化且$\beta_t=1$时，更新后$\bm S_t\bm k_t=\bm v_t$精确成立。

delta规则擅长精确修正，却没有快速清空大量过时内容的手段；衰减门恰好相反。Gated DeltaNet把两者结合\scite{nlp-llm-gateddelta2024}：
\begin{equation}
 \bm S_t=\alpha_t\bm S_{t-1}(\bm I-\beta_t\bm k_t\bm k_t^{\top})+\beta_t\bm v_t\bm k_t^{\top},\qquad
 \bm o_t=\bm S_t\bm q_t,
 \label{eq:nlp-llm-gated-delta}
\end{equation}
其中$\alpha_t\in(0,1)$是数据相关的标量衰减门。$\alpha_t\to0$时整段旧记忆被迅速清除，$\alpha_t\to1$时退化为纯delta规则，模型可以按内容在“整体遗忘”和“定点改写”之间切换。Kimi Linear中的Kimi Delta Attention把标量$\alpha_t$换成逐通道的对角门$\operatorname{Diag}(\bm\alpha_t)$，相当于把GLA的细粒度遗忘与delta规则结合\scite{nlp-llm-kimilinear2025}。

delta规则的状态转移$\bm I-\beta_t\bm k_t\bm k_t^{\top}$依赖状态之外的输入，且相邻步的转移矩阵不可交换，不能像衰减那样简单地写成累积乘积。Yang等指出这些矩阵的连乘可以写成“单位阵减去若干低秩外积之和”的WY表示，从而在块内用矩阵乘法计算，块间仍只传递状态，使delta类模型获得与GLA相当的分块并行训练\scite{nlp-llm-deltanet2024}。

\subsubsection{选择性状态空间模型：Mamba}
\label{sec:nlp-llm-mamba}

状态空间模型从另一条路线得到递归汇聚。\term{状态空间模型}（state space model, SSM）源自连续系统$\bm h'(\tau)=\bm A\bm h(\tau)+\bm Bx(\tau)$、$y(\tau)=\bm C\bm h(\tau)$，以步长$\Delta$离散化后得到一般的递推
\begin{equation}
 \bm h_t=\overline{\bm A}_t\bm h_{t-1}+\overline{\bm B}_tx_t,\qquad
 y_t=\bm C_t\bm h_t,
 \label{eq:nlp-llm-state-space}
\end{equation}
这里按单个通道书写，$x_t$是标量输入，$\bm h_t\in\mathbb R^{m}$是该通道的状态；零阶保持离散化给出$\overline{\bm A}_t=\exp(\Delta_t\bm A)$，实践中$\overline{\bm B}_t$常近似为$\Delta_t\bm B_t$。S4一类早期模型的参数不随输入变化，整段输出可以写成一次长卷积；但它们对每个位置都用同样的保留率，不能按内容决定记住什么。

Mamba让步长和读写向量都依赖当前输入\scite{gu2023mamba}：$\Delta_t=\operatorname{softplus}(\bm w_{\Delta}^{\top}\bm x_t+b_{\Delta})$，$\bm B_t=\bm W_B\bm x_t$，$\bm C_t=\bm W_C\bm x_t$，其中$\bm x_t$是该层所有通道的输入向量，$\bm A$为固定的对角负矩阵。$\Delta_t$起到门控作用：$\Delta_t$大时$\overline{\bm A}_t$接近0，状态被当前输入覆盖；$\Delta_t$小时状态几乎原样保留。输入相关的参数使整段响应不再是固定卷积核，Mamba改用并行扫描：仿射更新$\bm h\mapsto\overline{\bm A}_t\bm h+\overline{\bm B}_tx_t$的复合仍是仿射映射，且满足结合律，可以像前缀和一样在$O(\log T)$并行步内算出全部状态。选择性更新、完整模块的两条分支和硬件感知扫描的推导可参考\BookChapterRef{chap:recurrent-networks}{卷五的循环神经网络章}，这里只保留语言模型需要的部分。

Mamba-2进一步把每个头的$\bm A$限制为标量乘单位阵$a_t\bm I$。此时把递推展开，得
\[
 y_t=\sum_{j\le t}\Bigl(\prod_{k=j+1}^{t}a_k\Bigr)(\bm C_t^{\top}\bm B_j)\,x_j ,
\]
与式\eqref{eq:nlp-llm-retnet}比较可见：$\bm C_t$扮演查询，$\bm B_j$扮演键，$x_j$扮演值，$\prod a_k$是数据相关的标量衰减。也就是说，这一受限的状态空间模型与带标量门的线性注意力是同一计算的两种写法，Dao与Gu称之为结构化状态空间对偶\scite{dao2024mamba2}。这一对偶要求标量结构的转移矩阵，并不意味着所有状态空间模型都等价于softmax注意力。由此，线性注意力、RetNet、GLA、Mamba-2与Gated DeltaNet可以放进同一框架：都维护矩阵状态，区别在于转移矩阵是恒等、固定标量、数据相关对角阵，还是带擦除项的低秩修正。

\subsubsection{线性层与全注意力层的混合}
\label{sec:nlp-llm-hybrid-layers}

固定大小状态在精确回读上存在硬限制：若任务要求从十万词元前复制一串任意编号，状态容量再大也只能保留有限的信息。一种折中是让多数层使用线性注意力或状态空间层，按固定间隔插入少量全注意力层，由后者承担精确回读，前者承担廉价的局部与中程汇聚。Jamba按1层注意力配7层Mamba的比例交错堆叠，并每隔一层把前馈层换成混合专家层\scite{nlp-llm-jamba2024}；Kimi Linear按3层Kimi Delta Attention配1层MLA的比例混合\scite{nlp-llm-kimilinear2025}。

混合结构的缓存不能再称为固定大小。全注意力层的键值仍随长度线性增长，只是层数变少：若四分之一的层为全注意力，其余层为固定状态，则缓存约为同配置纯注意力模型的四分之一，再加上与长度无关的递归状态。训练计算同样只在少量层上保留平方项。混合比例、全注意力层放在哪些深度，以及全注意力层是否使用位置编码，都会影响长文本上的表现，需要在具体设置下实验确定。

\subsection{长上下文结构的比较}
\label{sec:nlp-llm-context-comparison}

表~\ref{tab:nlp-llm-context-comparison}把前面的方法放在同一组指标下比较。训练复杂度按一层、序列长度$T$、宽度$d$计；状态量按一层、生成到第$T$个词元时需要保存的元素数计；读取量指生成一个词元时需要访问的历史位置或状态。

\begin{table*}[htbp]
 \small
 \begin{tabularx}{\linewidth}{@{}p{17mm}p{27mm}p{30mm}p{23mm}p{21mm}X@{}}
  \toprule
  方法 & 训练复杂度 & 每层推断状态量 & 每步读取 & 精确回读 & 代表模型\\
  \midrule
  全注意力 & $O(T^2d)$ & $2Hd_hT$ & 全部$T$个位置 & 任意位置 & GPT-2\\
  GQA & $O(T^2d)$ & $2H_{\mathrm{kv}}d_hT$ & 全部$T$个位置 & 任意位置 & LLaMA-3\\
  MLA & $O(T^2d)$ & $(d_c+d_R)T$ & 全部$T$个位置 & 任意位置 & DeepSeek-V2/V3\\
  滑动窗口 & $O(Twd)$ & $2H_{\mathrm{kv}}d_hw$ & 最近$w$个位置 & 仅窗口内 & Mistral 7B\\
  NSA & $O(T(T/\delta+nl'+w)d)$ & 全部键值另加压缩键值 & $T/\delta+nl'+w$ & 仅被选中块 & 原文实验模型\\
  线性注意力 & $O(TCd+Td_kd_v)$ & $d_kd_v$（每头） & 一个状态矩阵 & 否 & 早期线性Transformer\\
  门控delta类 & 同上，另加WY计算 & $d_kd_v$（每头） & 一个状态矩阵 & 否，可定点改写 & Gated DeltaNet\\
  Mamba & $O(Td_{\mathrm{in}}m)$ & $d_{\mathrm{in}}m$ & 一个状态 & 否 & Mamba、Mamba-2\\
  混合结构 & 少数层$O(T^2d)$ & 少数层随$T$增长 & 少数层读全部 & 由全注意力层承担 & Jamba、Kimi Linear\\
  \bottomrule
 \end{tabularx}
 \caption{长上下文结构在训练计算、推断状态、读取量与回读能力上的比较。符号沿用正文：$C$为分块长度，$d_{\mathrm{in}}$与$m$为状态空间层的通道数和每通道状态宽度，NSA的$\delta,l',n$为压缩步长、选择块长与选块数。表中只比较资源与访问方式，不比较任务效果。}
 \label{tab:nlp-llm-context-comparison}
\end{table*}

这张表说明各方法节省的是不同的资源。键值压缩减少每个位置的表示量，读取范围和精确回读能力保持不变；稀疏读取减少每次访问的位置数，但缓存仍要完整保存，固定模式与按内容选择分别承担接力损失和选漏风险；递归汇聚把历史压成定长状态，换来与长度无关的推断代价，放弃的是逐位置精确回读；混合结构把精确回读集中到少数层。这些做法也可以组合，例如Kimi Linear同时使用了MLA与delta类层。缓存更小或复杂度更低并不能单独推出长文本任务表现更好，长上下文能力需要按第\ref{chap:nlp-evaluation}章的长上下文评测另行检查；把训练好的模型扩展到更长上下文的位置编码调整，则在第\ref{sec:nlp-llm-context-extension}节讨论。
